\section{Related Work}

\subsection{Filter Pruning Methods}
Model pruning reduces the computational cost of deep neural networks by removing redundant parameters. Structured pruning, particularly filter pruning, has gained popularity as it does not require specialized hardware support compared to unstructured pruning \citep{han2015deep}. Existing methods can be broadly categorized into magnitude-based, geometric, and reconstruction-based approaches.

\textbf{Magnitude-based Pruning.} Early works like L1-Norm pruning \citep{li2016pruning} and Network Slimming \citep{liu2017learning} assume that filters with smaller weights contribute less to the network's output. While simple and effective, this assumption does not always hold, as small-magnitude weights can still be critical for preserving high-level semantic information \citep{he2019filter}. Our GRA-CNN differs by evaluating the *relational importance* of channels across the entire dataset sequence, rather than relying on static weight magnitudes.

\textbf{Geometric and Rank-based Pruning.} To address the limitations of magnitude-based methods, geometric approaches like FPGM \citep{he2019filter} utilize the geometric median to identify redundant filters that can be replaced by others. Similarly, HRank \citep{lin2020hrank} prunes filters based on the rank of their feature maps, assuming low-rank features are less informative. While these methods capture structural redundancy, they often process each layer independently or rely on local statistical properties. In contrast, GRA-CNN introduces a global perspective by aligning channel importance with the decision boundary through semantic sequence analysis.

\textbf{Reconstruction and Global Methods.} Methods like ThiNet \citep{luo2017thinet} and CP \citep{he2017channel} focus on reconstructing the output of the next layer to guide pruning. More recent approaches like CHIP \citep{sui2021chip} and ResRep \citep{ding2021resrep} explore channel independence and re-parameterization. However, these methods typically require complex optimization procedures or extensive fine-tuning.

\subsection{Global Relation Analysis in Vision}
Gray Relational Analysis (GRA) has been widely used in multi-criteria decision making to quantify the correlation between a reference sequence and comparative sequences \citep{deng1982control}. In computer vision, relation modules have been explored to capture long-range dependencies, such as in Relation Networks \citep{hu2018squeeze} and Non-local Neural Networks \citep{wang2018non}. However, the application of GRA for quantifying channel importance in model compression remains unexplored. Unlike attention mechanisms that focus on spatial or channel-wise re-calibration during inference \citep{hu2018squeeze}, GRA-CNN utilizes GRA as a static importance metric for pruning, bridging the gap between classic system analysis and modern deep learning compression.

\subsection{Recent Advances (2022-2024)}
Recent literature has shifted towards dynamic and semantic-aware pruning. For instance, dynamic pruning methods \citep{lin2023dynamic} adjust network structure during inference based on input difficulty, but introduce runtime overhead. Semantic preservation approaches \citep{wang2023semantic} attempt to align pruned feature maps with pre-trained teachers using distillation. In comparison, GRA-CNN offers a static, inference-efficient solution that implicitly preserves semantic alignment through its logit-aware reference sequence, achieving state-of-the-art trade-offs without the need for complex distillation pipelines or dynamic routing.
